\documentclass[12pt,a4paper]{amsart}

\usepackage{amsmath,amssymb,amsthm}
\usepackage{mathtools}
\usepackage[utf8]{inputenc}
\usepackage[T1]{fontenc}
\usepackage{hyperref}
\usepackage{enumitem,booktabs,array}
\usepackage{geometry}
\geometry{margin=2.5cm}

% -------------------------------------------------------
% Theorem environments
% -------------------------------------------------------
\newtheorem{theorem}{Theorem}[section]
\newtheorem{lemma}[theorem]{Lemma}
\newtheorem{corollary}[theorem]{Corollary}
\newtheorem{proposition}[theorem]{Proposition}
\theoremstyle{definition}
\newtheorem{definition}[theorem]{Definition}
\newtheorem{example}[theorem]{Example}
\newtheorem{remark}[theorem]{Remark}
\newtheorem{conjecture}[theorem]{Conjecture}

% -------------------------------------------------------
% Custom commands
% -------------------------------------------------------
\newcommand{\N}{\mathbb{N}}
\newcommand{\Z}{\mathbb{Z}}
\newcommand{\Q}{\mathbb{Q}}
\newcommand{\R}{\mathbb{R}}
\newcommand{\C}{\mathbb{C}}
\newcommand{\F}{\mathbb{F}}
\newcommand{\Zp}{\mathbb{Z}_p}
\newcommand{\Qp}{\mathbb{Q}_p}
\newcommand{\A}{\mathbb{A}}
\newcommand{\GL}{\mathrm{GL}}
\newcommand{\SL}{\mathrm{SL}}
\newcommand{\GQ}{\mathrm{Gal}(\bar{\Q}/\Q)}
\newcommand{\Gal}{\mathrm{Gal}}
\newcommand{\Frob}{\mathrm{Frob}}
\newcommand{\WD}{\mathrm{WD}}
\newcommand{\WF}{\mathrm{WF}}
\newcommand{\rec}{\mathrm{rec}}
\newcommand{\Ind}{\mathrm{Ind}}
\newcommand{\Art}{\mathrm{Art}}
\DeclareMathOperator{\ord}{ord}
\DeclareMathOperator{\Hom}{Hom}
\DeclareMathOperator{\End}{End}
\DeclareMathOperator{\Aut}{Aut}
\DeclareMathOperator{\tr}{tr}
\DeclareMathOperator{\nr}{N}

% -------------------------------------------------------
\begin{document}
% -------------------------------------------------------

\title{The Langlands Programme:\\
       $L$-functions, Automorphic Representations,\\
       and the Global Correspondence}

\author{Michael Fuhrmann}
\address{Specialist Mathematics Project, Build 120}
\email{specialist-maths@localhost}
\date{\today}

\begin{abstract}
The Langlands programme, initiated by Robert Langlands in a 1967 letter to
André Weil, proposes a vast web of conjectures relating number theory,
automorphic forms, and representation theory.  At its heart is the
\emph{Langlands correspondence}: a (conjectural) bijection between
$n$-dimensional representations of the absolute Galois group
$\GQ$ (or more generally of the Weil group) and automorphic representations
of $\GL_n(\mathbb{A}_\Q)$.  This paper develops the necessary background
(Galois representations, $L$-functions, automorphic forms, the local
Langlands correspondence) and presents the major known cases and applications,
including Wiles' proof of Fermat's Last Theorem as a landmark instance of
the programme.
\end{abstract}

\maketitle
\tableofcontents

% ================================================================
\section{Motivation: From Quadratic Reciprocity to Langlands}
% ================================================================

The Langlands programme generalises a long chain of reciprocity laws:

\begin{itemize}
  \item \textbf{Quadratic reciprocity} (Gauss 1796): For distinct odd primes $p, q$,
    $\bigl(\frac{p}{q}\bigr)\bigl(\frac{q}{p}\bigr) = (-1)^{\frac{p-1}{2}\cdot\frac{q-1}{2}}$.
  \item \textbf{Class field theory} (Artin 1927): Every abelian extension $L/K$
    corresponds to an open subgroup of the idèle class group $\mathbb{A}_K^\times/K^\times$.
    The Artin map $\Art_{L/K}$ identifies Galois groups with quotients of this group.
  \item \textbf{Langlands}: Non-abelian extensions correspond to automorphic
    representations of $\GL_n$ rather than of $\GL_1$.
\end{itemize}

The key insight: the arithmetic of number fields (Galois groups) and the
analysis of automorphic forms (harmonic analysis on adèle groups) are
secretly the same subject.

% ================================================================
\section{Galois Representations}
% ================================================================

\begin{definition}[Galois representation]
\label{def:galois_rep}
A \emph{(continuous) $\ell$-adic Galois representation} of degree $n$ is a
continuous group homomorphism
\[
  \rho : \GQ \longrightarrow \GL_n(\bar{\Q}_\ell)
\]
(or into $\GL_n(\Zp)$, $\GL_n(\Qp)$, etc.) where $\bar{\Q}_\ell$ carries
the $\ell$-adic topology.  Such $\rho$ is \emph{unramified at $p$} if
$\rho|_{I_p} = 1$ for the inertia group $I_p \subset \GQ$.
\end{definition}

\begin{example}[Cyclotomic character]
The \emph{cyclotomic character} $\chi_\ell : \GQ \to \Z_\ell^\times$ is defined by
$\sigma(\zeta_{\ell^n}) = \zeta_{\ell^n}^{\chi_\ell(\sigma)}$ for all $\ell^n$-th
roots of unity.  This is a 1-dimensional representation.
\end{example}

\begin{example}[Tate module of an elliptic curve]
Let $E/\Q$ be an elliptic curve.  The \emph{$\ell$-adic Tate module} is
\[
  T_\ell(E) = \varprojlim_n E[\ell^n](\bar\Q) \;\cong\; \Z_\ell^2.
\]
The Galois action on $\ell$-power torsion gives a 2-dimensional representation
\[
  \rho_{E,\ell} : \GQ \longrightarrow \GL_2(\Z_\ell) \hookrightarrow \GL_2(\Q_\ell).
\]
For a prime $p \nmid \ell\cdot\Delta(E)$ (good reduction), $\rho_{E,\ell}$
is unramified at $p$ and
\[
  \tr(\rho_{E,\ell}(\Frob_p)) = a_p(E) = p + 1 - |E(\F_p)|.
\]
\end{example}

\begin{definition}[$L$-function of a Galois representation]
\label{def:galois_l}
For $\rho : \GQ \to \GL_n(\Q_\ell)$ unramified at all but finitely many primes,
the \emph{$L$-function} is
\[
  L(s, \rho) = \prod_p \det\!\bigl(1 - \rho(\Frob_p)\,p^{-s}\bigr)^{-1},
\]
where the product runs over primes $p$ where $\rho$ is unramified.
At ramified primes, one uses the \emph{local Euler factor}
$L(s, \rho|_{W_{\Q_p}})$ from the Weil--Deligne group.
\end{definition}

\begin{example}[$L$-function of an elliptic curve]
For $E/\Q$: $L(s, E) = L(s, \rho_{E,\ell})$.  By Wiles' theorem (Shimura--Taniyama),
this equals the $L$-function of a weight-$2$ cusp form $f_E$.  The Riemann
hypothesis for $L(s,E)$ (i.e.\ zeroes on $\Re(s) = 1$) follows from
Deligne's theorem on the Ramanujan conjecture.
\end{example}

% ================================================================
\section{Automorphic Representations}
% ================================================================

\begin{definition}[Adèle ring]
\label{def:adeles}
The \emph{adèle ring} of $\Q$ is the restricted product
\[
  \mathbb{A} = \mathbb{A}_\Q = \R \times \prod_p{}' \Q_p
  = \bigl\{(x_\infty, x_2, x_3, x_5, \ldots) : x_p \in \Z_p \text{ for almost all } p\bigr\}.
\]
It is a locally compact ring.  The diagonal embedding $\Q \hookrightarrow \mathbb{A}$
(as constant sequences) makes $\Q$ a discrete cocompact subring.
\end{definition}

\begin{definition}[Automorphic representation]
\label{def:automorphic_rep}
An \emph{automorphic representation} of $\GL_n(\mathbb{A})$ is an irreducible
admissible representation $\pi = \otimes'_v \pi_v$ (restricted tensor product over
all places $v$: the primes $p$ and $v = \infty$) that occurs in the space
$L^2(\GL_n(\Q) \backslash \GL_n(\mathbb{A}), \omega)$ for some central character $\omega$.

A \emph{cuspidal} automorphic representation is one for which all constant
terms along proper parabolic subgroups vanish.
\end{definition}

\begin{remark}
The decomposition $\pi = \otimes'_v \pi_v$ (tensor product of local representations)
is a fundamental structural theorem (Flath's theorem).  Each local factor $\pi_p$
is an irreducible admissible representation of $\GL_n(\Q_p)$.
\end{remark}

\begin{example}[Classical modular forms as automorphic representations]
A normalised Hecke eigenform $f \in S_k(\Gamma_0(N), \chi)$ (weight $k$,
level $N$, nebentypus $\chi$) corresponds to a cuspidal automorphic
representation $\pi_f$ of $\GL_2(\mathbb{A})$ via
\[
  f(z) \leftrightarrow \phi_f : g \mapsto f\!\left(\frac{ai+b}{ci+d}\right)
  (ci+d)^{-k}, \quad g = \begin{pmatrix} a & b \\ c & d \end{pmatrix}.
\]
The $L$-function of $\pi_f$ is $L(s, f) = \sum_{n \ge 1} a_n n^{-s}$.
\end{example}

\begin{definition}[$L$-function of an automorphic representation]
\label{def:automorphic_l}
The \emph{(standard) $L$-function} of $\pi = \otimes_v \pi_v$ is
\[
  L(s, \pi) = \prod_v L(s, \pi_v),
\]
where $L(s, \pi_p) = \det(1 - A_p p^{-s})^{-1}$ (for unramified $\pi_p$,
with $A_p \in \GL_n(\C)$ the Satake parameter).
These $L$-functions are known (by work of Godement--Jacquet) to have analytic
continuation and functional equation.
\end{definition}

% ================================================================
\section{The Local Langlands Correspondence}
% ================================================================

\begin{definition}[Weil--Deligne group]
\label{def:weil_deligne}
For a $p$-adic field $F$ (e.g.\ $F = \Q_p$), the \emph{Weil group} $W_F$ is
the dense subgroup of $\Gal(\bar{F}/F)$ consisting of elements whose image
in $\Gal(\bar{\F}_p/\F_p)$ is an integral power of Frobenius.
The \emph{Weil--Deligne group} $\WD(F)$ is the semi-direct product
$W_F \ltimes \C$ with $w \cdot n \cdot w^{-1} = \|w\|\cdot n$
(where $\|w\| = q^{-v(w)}$ for the Frobenius valuation $v$).
\end{definition}

\begin{theorem}[Local Langlands for $\GL_n$]
\label{thm:local_langlands}
Let $F$ be a finite extension of $\Q_p$.  There is a canonical bijection
\[
  \mathrm{LLC}_F :
  \left\{\parbox{5cm}{\centering $n$-dim.\ irreducible representations of $\WD(F)$}\right\}
  \;\longleftrightarrow\;
  \left\{\parbox{5cm}{\centering irreducible admissible representations of $\GL_n(F)$}\right\}
\]
satisfying:
\begin{enumerate}[label=(\roman*)]
  \item $L(s, \pi) = L(s, \mathrm{LLC}_F(\pi))$ (compatibility of $L$-factors),
  \item $\varepsilon(s, \pi, \psi) = \varepsilon(s, \mathrm{LLC}_F(\pi), \psi)$ (compatibility of $\varepsilon$-factors),
  \item $\mathrm{LLC}_F(\chi \circ \det \otimes \pi) = (\chi \circ \Art_F) \otimes \mathrm{LLC}_F(\pi)$.
\end{enumerate}
\end{theorem}

\begin{proof}[Status]
For $\GL_1(F) = F^\times$: Class field theory (Artin map $\Art_F : F^\times \to W_F^{ab}$). \\
For $\mathrm{GL}_2(F)$: Proved by Kutzko (1980). \\
For $\mathrm{GL}_n(F)$: Proved independently by Henniart (2000) and Harris--Taylor (2001).
\end{proof}

% ================================================================
\section{The Global Langlands Correspondence}
% ================================================================

\begin{conjecture}[Global Langlands, $n = 1$ case]
\label{conj:gl1}
For $n = 1$: Artin's abelian class field theory gives a bijection between
continuous characters $\chi : \GQ \to \C^\times$ (= 1-dim.\ Galois reps.)
and Hecke characters (automorphic forms of $\GL_1(\A)$):
\[
  \chi \leftrightarrow \omega_\chi \in \mathrm{Aut}(\GL_1(\A)).
\]
This is the abelian case — fully proved by class field theory.
\end{conjecture}

\begin{conjecture}[Global Langlands for $\GL_n$]
\label{conj:global_langlands}
There is a canonical bijection (the \emph{Langlands correspondence})
\[
  \mathcal{L} :
  \left\{\parbox{5.5cm}{\centering irreducible $n$-dim.\ $\ell$-adic representations $\rho$ of $\GQ$
  (semisimple, geometric)}\right\}
  \;\longleftrightarrow\;
  \left\{\parbox{5.5cm}{\centering cuspidal automorphic representations
  $\pi$ of $\GL_n(\mathbb{A})$}\right\}
\]
such that $L(s, \rho) = L(s, \pi)$ (equality of $L$-functions).
\end{conjecture}

\begin{remark}
The correspondence is expected to preserve:
\begin{itemize}
  \item Local $L$-factors: $L(s, \rho_p) = L(s, \pi_p)$ for all primes $p$.
  \item Functional equation: both $L$-functions satisfy the same $\Gamma$-factor.
  \item Local Langlands at each place.
\end{itemize}
\end{remark}

\begin{theorem}[Langlands--Tunnell, $n = 2$, solvable image]
\label{thm:langlands_tunnell}
If $\rho : \GQ \to \GL_2(\C)$ is continuous with solvable image, then
$L(s, \rho)$ is the $L$-function of an automorphic form on $\GL_2(\A)$.
\end{theorem}

This theorem was crucial in Wiles' proof of Fermat's Last Theorem.

% ================================================================
\section{Shimura--Taniyama--Wiles: An Explicit Case}
% ================================================================

\begin{theorem}[Shimura--Taniyama--Wiles--Breuil--Conrad--Diamond--Taylor]
\label{thm:stw}
Every elliptic curve $E/\Q$ is \emph{modular}: there exists a Hecke eigenform
$f_E \in S_2(\Gamma_0(N))$ (for the conductor $N$ of $E$) such that
\[
  L(s, E) = L(s, f_E).
\]
Equivalently, the 2-dimensional Galois representation $\rho_{E,\ell}$ corresponds
to the automorphic representation $\pi_{f_E}$ under the Langlands correspondence.
\end{theorem}

\begin{proof}[History and proof sketch]
\textbf{Strategy (Wiles 1995):}
\begin{enumerate}[label=\arabic*.]
  \item \emph{Residual representation:} The mod-$\ell$ representation
    $\bar\rho_{E,\ell} : \GQ \to \GL_2(\F_\ell)$ is modular by Langlands--Tunnell
    (for $\ell = 3$, using solvability of $\GL_2(\F_3)$).
  \item \emph{Deformation theory:} Mazur's deformation theory parametrises
    liftings of $\bar\rho$ by a universal deformation ring $R$.  Wiles proves
    $R \cong T$ where $T$ is the Hecke algebra — the \emph{$R = T$ theorem}.
  \item \emph{$R = T$:} Using a \emph{numerical criterion} (Wiles) and the
    Horizontal Iwasawa theory, the surjection $R \twoheadrightarrow T$ is shown
    to be an isomorphism.
  \item \emph{Conclusion:} Every $E$ with semistable reduction was handled by
    Wiles; Taylor--Wiles completed the semistable case; BCDT (2001) completed
    the general case.
\end{enumerate}
\end{proof}

\begin{corollary}[Fermat's Last Theorem (Wiles 1995)]
The equation $x^n + y^n = z^n$ has no solution in positive integers for $n \ge 3$.
\end{corollary}

\begin{proof}[Proof sketch]
Ribet (1990) proved: if a solution $(a, b, c)$ exists, the \emph{Frey curve}
$E: y^2 = x(x - a^n)(x + b^n)$ would be a non-modular elliptic curve
(contradicting the Shimura--Taniyama conjecture).  Since $E$ is modular
(Wiles' theorem), no solution exists.
\end{proof}

% ================================================================
\section{The Fontaine--Mazur Conjecture}
% ================================================================

\begin{conjecture}[Fontaine--Mazur, 1995]
\label{conj:fontaine_mazur}
An irreducible $n$-dimensional $p$-adic representation $\rho : \GQ \to \GL_n(\Q_p)$
comes from algebraic geometry (i.e.\ is a sub-representation of the étale
cohomology of a smooth projective variety over $\Q$) if and only if it is:
\begin{enumerate}[label=(\roman*)]
  \item \emph{Geometric}: unramified outside a finite set of primes and de Rham
    at $p$ (in the sense of $p$-adic Hodge theory),
  \item \emph{Odd}: the complex conjugation acts with eigenvalues $\pm 1$ (for $n=2$).
\end{enumerate}
Such geometric representations should correspond to cuspidal automorphic representations
(under the Langlands correspondence).
\end{conjecture}

\begin{remark}
This conjecture is the precise form of the global Langlands programme for
$p$-adic representations.  Wiles' theorem proves a special case
(weight-$2$ representations from elliptic curves).  The $p$-adic Langlands
programme (Colmez, Breuil, Emerton, Kisin) proves further cases.
\end{remark}

% ================================================================
\section{The Langlands Functoriality Conjecture}
% ================================================================

\begin{conjecture}[Langlands Functoriality]
\label{conj:functoriality}
Let $G$ and $H$ be reductive groups over $\Q$, and ${}^L G$, ${}^L H$ be
their \emph{Langlands dual groups}.  For any $L$-group homomorphism
$\phi : {}^L H \to {}^L G$, there exists a natural \emph{transfer map}
(or \emph{functorial lift}) on automorphic representations:
\[
  \text{Aut}(H) \ni \pi \;\longmapsto\; \phi_* \pi \in \text{Aut}(G),
\]
preserving $L$-functions: $L(s, \pi, r \circ \phi) = L(s, \phi_*\pi, r)$
for any finite-dimensional representation $r$ of ${}^L G$.
\end{conjecture}

\begin{example}[Base change for $\GL_n$]
The \emph{base change lift} $\pi \mapsto \mathrm{BC}_{L/K}(\pi)$ for cyclic
extensions $L/K$ was proved by Langlands ($n=2$, 1980) and Arthur--Clozel ($n$ general, 1989).
This gives: if $\pi$ is an automorphic representation of $\GL_n(\mathbb{A}_K)$,
then there exists a functorial lift to $\GL_n(\mathbb{A}_L)$.
\end{example}

\begin{example}[Symmetric power lifts]
The $m$-th symmetric power lift $\text{Sym}^m : \GL_2 \to \GL_{m+1}$ would
map $\pi_f$ (modular form of weight 2) to an automorphic representation of
$\GL_{m+1}$.  Known: $m = 2, 3, 4$ (Kim--Shahidi 2002).  Open for $m \ge 5$.
\end{example}

% ================================================================
\section{The $p$-adic Langlands Programme}
% ================================================================

\begin{remark}[Motivation]
The classical Langlands correspondence works over $\C$ (or $\bar\Q_\ell$).
The \emph{$p$-adic Langlands programme} seeks to replace $\GL_n(\Q_\ell)$ by
$\GL_n(\Q_p)$ and complex representations by $p$-adic ones, motivated by:
\begin{itemize}
  \item Kisin's work on potentially semi-stable deformation rings.
  \item Colmez's $p$-adic Local Langlands for $\GL_2(\Q_p)$ (2010).
  \item Emerton's completed cohomology approach.
\end{itemize}
\end{remark}

\begin{theorem}[$p$-adic Local Langlands for $\GL_2(\Q_p)$, Colmez 2010]
\label{thm:colmez}
There is a canonical bijection (the \emph{$p$-adic local Langlands correspondence})
\[
  \mathrm{LLC}_p^{\mathrm{p-adic}} :
  \left\{\parbox{5cm}{\centering 2-dim.\ irreducible $p$-adic representations
  of $\Gal(\bar\Q_p/\Q_p)$}\right\}
  \;\longleftrightarrow\;
  \left\{\parbox{5cm}{\centering admissible unitary Banach space representations of $\GL_2(\Q_p)$}\right\}
\]
compatible with the classical local Langlands via reduction mod $p$.
\end{theorem}

% ================================================================
\section{Known Cases and Open Problems}
% ================================================================

\begin{remark}[Summary of known cases]
\begin{center}
\begin{tabular}{l|l|l}
\textbf{Case} & \textbf{Groups} & \textbf{Status} \\
\hline
$n = 1$ (abelian) & $\GL_1$ & Proved (class field theory) \\
$n = 2$, solvable image & $\GL_2$ & Proved (Langlands--Tunnell) \\
Elliptic curves over $\Q$ & $\GL_2$ & Proved (Wiles + BCDT) \\
Local LLC for $\GL_n$ & $\GL_n(F)$, $F/\Q_p$ & Proved (Harris--Taylor, Henniart) \\
Base change $\GL_n$ & $\GL_n$ & Proved (Arthur--Clozel) \\
$\text{Sym}^m$ for $m \le 4$ & $\GL_2 \to \GL_{m+1}$ & Proved (Kim--Shahidi) \\
$p$-adic LLC for $\GL_2(\Q_p)$ & $\GL_2$ & Proved (Colmez) \\
\hline
Global LLC for $\GL_n$ ($n \ge 2$) & $\GL_n$ & \textbf{OPEN} \\
Fontaine--Mazur (general) & $\GL_n$ & Partial (Wiles $n=2$, $k=2$) \\
Non-abelian Artin conjecture & general $G$ & \textbf{OPEN} \\
Geometric Langlands & curves over $\C$ & Recent progress (2024) \\
\end{tabular}
\end{center}
\end{remark}

\begin{remark}[Geometric Langlands 2024]
In 2024, a team led by Gaitsgory completed the proof of the
\emph{geometric Langlands conjecture} for curves over algebraically closed
fields of characteristic 0, in the de Rham setting.  This is a major advance
in the programme, though the arithmetic (number field) case remains open.
\end{remark}

% ================================================================
\section{Summary}
% ================================================================

\begin{remark}[The Big Picture]
The Langlands programme creates a dictionary:
\begin{center}
\begin{tabular}{l|l}
\textbf{Number Theory / Geometry} & \textbf{Automorphic Forms / Representation Theory} \\
\hline
Galois group $\GQ$ & Algebraic group $G$ over $\A$ \\
$n$-dim.\ Galois representation & Automorphic representation of $\GL_n(\A)$ \\
Zeta/Hasse--Weil $L$-function & Automorphic $L$-function \\
Ramification at $p$ & Local representation $\pi_p$ \\
Frobenius at $p$ & Hecke eigenvalue $\lambda_p$ \\
Fermat curve / elliptic curve & Weight-2 cusp form \\
\end{tabular}
\end{center}
The unification of these two worlds — arithmetic and analytic — is the
deepest theme in modern mathematics.
\end{remark}

\begin{thebibliography}{99}

\bibitem{Langlands1970}
R.~P.~Langlands,
\textit{Problems in the theory of automorphic forms},
in: Lectures in Modern Analysis and Applications III,
Lecture Notes in Math.\ \textbf{170}, Springer, 1970, 18--61.

\bibitem{HarrisTaylor2001}
M.~Harris and R.~Taylor,
\textit{The geometry and cohomology of some simple Shimura varieties},
Ann.\ of Math.\ Studies \textbf{151}, Princeton Univ.\ Press, 2001.

\bibitem{Henniart2000}
G.~Henniart,
\textit{Une preuve simple des conjectures de Langlands pour $\GL(n)$ sur un corps $p$-adique},
Invent.\ Math.\ \textbf{139} (2000), 439--455.

\bibitem{Wiles1995}
A.~Wiles,
\textit{Modular elliptic curves and Fermat's Last Theorem},
Ann.\ of Math.\ (2) \textbf{141} (1995), 443--551.

\bibitem{BCDT2001}
C.~Breuil, B.~Conrad, F.~Diamond, R.~Taylor,
\textit{On the modularity of elliptic curves over $\Q$: wild 3-adic exercises},
J.\ Amer.\ Math.\ Soc.\ \textbf{14} (2001), 843--939.

\bibitem{ArthurClozel1989}
J.~Arthur and L.~Clozel,
\textit{Simple Algebras, Base Change, and the Advanced Theory of the Trace Formula},
Ann.\ of Math.\ Studies \textbf{120}, Princeton Univ.\ Press, 1989.

\bibitem{Colmez2010}
P.~Colmez,
\textit{Repr\'esentations de $\GL_2(\Q_p)$ et $(\varphi,\Gamma)$-modules},
Ast\'erisque \textbf{330} (2010), 281--509.

\bibitem{Ribet1990}
K.~Ribet,
\textit{On modular representations of $\Gal(\bar\Q/\Q)$ arising from modular forms},
Invent.\ Math.\ \textbf{100} (1990), 431--476.

\bibitem{FrenkelBenZvi}
E.~Frenkel,
\textit{Langlands Correspondence for Loop Groups},
Cambridge Univ.\ Press, 2007.

\end{thebibliography}

\end{document}
